\documentclass[11pt]{article}

\usepackage[letterpaper,margin=0.92in]{geometry}
\usepackage{amsmath,amssymb,amsthm,mathtools,mathrsfs}
\usepackage{graphicx}
\usepackage{xcolor}
\usepackage{enumitem}
\usepackage{microtype}
\usepackage[hidelinks]{hyperref}

\newtheorem{theorem}{Theorem}
\newtheorem{lemma}[theorem]{Lemma}
\newtheorem{remark}[theorem]{Remark}
\newtheorem{definition}[theorem]{Definition}

\newcommand{\BH}{B(H)}
\newcommand{\KH}{K(H)}
\newcommand{\cc}{\mathfrak c}
\newcommand{\Istar}{\mathcal I_{\!*}}
\newcommand{\diag}{\operatorname{diag}}
\newcommand{\amp}{\mathsf D}

\setlist[itemize]{leftmargin=1.5em}
\setlist[enumerate]{leftmargin=1.7em}
\setlength{\parindent}{0pt}
\setlength{\parskip}{0.55em}

\title{\textbf{Zero Quotient Dimension Can Require Continuum Many Subideal Generators}\\
\large A ZFC counterexample to Question 5 of \emph{Subideals of Operators II}}
\author{Automated mathematical discovery run \texttt{fa\_banach\_001}\\
Agent \texttt{agent\_lane\_14}\qquad Model GPT5.6}
\date{23 July 2026}

\begin{document}
\maketitle

\begin{center}
\fcolorbox{blue!60!black}{blue!4}{
\begin{minipage}{0.91\textwidth}
\textbf{Status: candidate counterexample, likely valid, pending human review.}
The construction gives a proper \(K(H)\)-ideal with zero quotient dimension
but generator number exactly the continuum.  It refutes both sentences of
the source's Question 5 under their literal hypotheses.
\end{minipage}}
\end{center}

\section{Result at a glance}

Let \(H\) be a separable infinite-dimensional complex Hilbert space.
Patnaik and Weiss define, for a \(J\)-ideal \(\mathcal I\),
\[
 \mathcal I^0
   =\operatorname{span}(\mathcal I J+J\mathcal I)+J(\mathcal I)J,
\]
where \((\mathcal I)\) is the \(\BH\)-ideal generated by
\(\mathcal I\) \cite{PW12}.  They ask whether
\(\dim_{\mathbb C}(\mathcal I/\mathcal I^0)\) is the cardinality of a
generating set, or at least whether dimension below
\(\cc=2^{\aleph_0}\) forces a generating set of cardinality below \(\cc\).

We construct a proper \(\BH\)-ideal
\[
                 \Istar\subsetneq\KH
\]
such that
\[
 \Istar^2=\Istar,
 \qquad
 \min\{|S|:(S)_{\KH}=\Istar\}=\cc.
\]
Taking the ambient ideal \(J=\KH\) and \(\mathcal I=\Istar\),
idempotence gives
\[
 \mathcal I=\mathcal I^2
 \subset \mathcal I J
 \subset \mathcal I,
\]
so \(\mathcal I^0=\mathcal I\).  Therefore
\[
 \boxed{\dim_{\mathbb C}(\mathcal I/\mathcal I^0)=0<\cc
 \quad\text{while}\quad
 \min\{|S|:(S)_J=\mathcal I\}=\cc.}
\]

\clearpage
\section{Source question}

The questions appear in Section 6, page 9 of arXiv:1209.6323
\cite{PW12}.  The following is a full-width crop of the source page.

\begin{center}
  \centering
  \includegraphics[width=\textwidth]{figures/open_problem_crop.png}

  \smallskip
  \textit{Figure 1. The source's five questions.  This packet answers
  Question 5.}
\end{center}

\textbf{Question 5 (transcription).}
``Is the dimension of \(\mathcal I/\mathcal I^0\) the cardinality of some
generating set for \(\mathcal I\)?  At least if
\(\dim\mathcal I/\mathcal I^0<\cc\), must \(\mathcal I\) possess a
generating set of cardinality less than \(\cc\)?''

\section{Definitions and a finite-domination lemma}

For a subset \(S\) of a \(\BH\)-ideal \(J\), write \((S)\) for the
\(\BH\)-ideal generated by \(S\), and \((S)_J\) for the two-sided ideal of
the nonunital ring \(J\) generated by \(S\).  Let \(c_0^*\) be the cone of
positive nonincreasing null sequences.  For \(\xi\in c_0^*\), put
\[
 D_\xi=\diag(\xi_1,\xi_2,\ldots).
\]
For \(m\geq1\), the \(m\)-fold ampliation of a sequence is
\[
             (\amp_m\xi)_n=\xi_{\lceil n/m\rceil}.
\]

The proof uses the following elementary form of the characteristic-set
criterion.  We include the argument to make the packet independent of a
classification theorem.

\begin{lemma}[Finite ideal domination]\label{lem:finite-domination}
Let \(\xi,\eta_1,\ldots,\eta_q\in c_0^*\).  If
\[
             D_\xi\in(D_{\eta_1},\ldots,D_{\eta_q}),
\]
then there are \(C>0\) and positive integers \(m_i\) such that
\[
             \xi\leq C\sum_{i=1}^q\amp_{m_i}\eta_i
             \quad\text{coordinatewise}.
\]
\end{lemma}

\begin{proof}
Membership in the algebraically generated ideal expresses
\(D_\xi\) as a finite sum
\[
             D_\xi=\sum_{j=1}^r A_jD_{\eta_{i_j}}B_j,
             \qquad A_j,B_j\in\BH.
\]
For \(T_j=A_jD_{\eta_{i_j}}B_j\),
\[
 s_n(T_j)\leq\|A_j\|\,\|B_j\|\,(\eta_{i_j})_n.
\]
Put \(a=\lceil n/r\rceil\).  Then \(ra-r+1\leq n\), and the repeated
Ky Fan singular-number inequality \cite[Chapter 2]{Simon} gives
\[
\begin{aligned}
 \xi_n=s_n(D_\xi)
 &\leq s_{ra-r+1}\!\left(\sum_{j=1}^rT_j\right) \\
 &\leq\sum_{j=1}^rs_a(T_j)
 \leq\sum_{j=1}^r
       \|A_j\|\,\|B_j\|\,(\eta_{i_j})_{\lceil n/r\rceil}.
\end{aligned}
\]
Grouping equal \(\eta_i\)'s and enlarging one constant yields the stated
finite ampliation bound.
\end{proof}

\section{Proof intuition}

The quotient \(\mathcal I/\mathcal I^0\) sees generators only modulo products
with the ambient ideal.  An idempotent ideal can therefore be completely
invisible to this quotient: if \(\mathcal I^2=\mathcal I\), then
\(\mathcal I\KH=\mathcal I\) and the quotient is zero.

The difficulty is to make such an ideal require \emph{continuum many}
generators in ZFC.  We encode a continuum independent family of subsets of
\(\mathbb N\) into compact diagonal operators.  At infinitely many enormous
blocks, one selected sequence is relatively large while any prescribed
finite collection of the other sequences is extremely small.  Lemma
\ref{lem:finite-domination} then prevents the selected diagonal from lying in
the ideal generated by those competitors.

Finally, we include every successive square root of every diagonal generator.
Each generator is then literally the square of another generator, so the
resulting ideal is algebraically idempotent.  The independent-family
separation survives all these roots and forces the generator number to be
\(\cc\).

\section{The counterexample}

\subsection{A continuum independent family}

\begin{lemma}\label{lem:independent}
There is a family
\[
       \mathscr A=\{A_\alpha:\alpha<\cc\}\subset\mathcal P(\mathbb N)
\]
such that every finite Boolean cell is infinite.  Explicitly, for distinct
\(\alpha_0,\ldots,\alpha_q\) and arbitrary
\(\varepsilon_i\in\{0,1\}\), the set
\[
 \{k:\mathbf 1_{A_{\alpha_i}}(k)=\varepsilon_i
          \text{ for }0\leq i\leq q\}
\]
is infinite.
\end{lemma}

\begin{proof}
Let
\[
 \Omega=\bigcup_{n\geq1}\bigl(\{n\}\times\mathcal P(2^n)\bigr),
\]
where \(2^n\) denotes the finite set of binary strings of length \(n\).
The set \(\Omega\) is countable.  For \(x\in2^{\mathbb N}\), define
\[
 A_x=\{(n,F)\in\Omega:x|_n\in F\}.
\]
Given finitely many distinct \(x_i\), their length-\(n\) prefixes are
distinct for every sufficiently large \(n\).  For a prescribed Boolean
pattern, let \(F\subset2^n\) contain exactly those prefixes assigned the
value \(1\).  The point \((n,F)\) realizes the pattern for every sufficiently
large \(n\).  Identify \(\Omega\) with \(\mathbb N\).
\end{proof}

\subsection{Strongly incomparable null sequences}

Set
\[
 N_0=L_0=1,
 \qquad
 N_k=L_k=2^{2^k}\quad(k\geq1).
\]
For \(A\in\mathscr A\), define
\[
 r_A(k)=\max\bigl(\{0\}\cup(A\cap\{1,\ldots,k\})\bigr).
\]
On the block \(N_k\leq n<N_{k+1}\), put
\[
       f_A(n)=L_{r_A(k)},
       \qquad
       \xi_A(n)=e^{-f_A(n)}.
\]
The values before the first nontrivial block are set equal to \(e^{-L_0}\).
Since \(A\) is infinite, \(f_A\) is nondecreasing and tends to infinity.
Thus \(\xi_A\in c_0^*\).

\begin{lemma}[Strong separation]\label{lem:separation}
For each \(A\in\mathscr A\),
\[
 D_{\xi_A}\notin
 \bigl(D_{\xi_B^t}:B\in\mathscr A\setminus\{A\},\ t>0\bigr).
\]
\end{lemma}

\begin{proof}
Ideal membership uses only finitely many displayed generators.  If the
conclusion failed, Lemma \ref{lem:finite-domination} would give distinct
\(B_1,\ldots,B_q\neq A\), positive \(t_i\), integers \(m_i\), and \(C>0\)
such that
\[
        \xi_A\leq C\sum_{i=1}^q
                    \amp_{m_i}(\xi_{B_i}^{t_i}).
\tag{1}\label{eq:domination}
\]
Let \(M=\max_i m_i\) and \(t=\min_i t_i>0\).  Lemma
\ref{lem:independent} gives infinitely many \(k\) for which
\[
              k\notin A,\qquad k\in B_i\quad(1\leq i\leq q).
\]
Choose such a sufficiently large \(k\), and set \(n=MN_k\).  Since
\(N_{k+1}/N_k\to\infty\), the integers \(n\) and
\(\lceil n/m_i\rceil\) all lie in the \(k\)-th block.  Consequently
\[
        f_A(n)\leq L_{k-1},
        \qquad
        f_{B_i}(\lceil n/m_i\rceil)=L_k.
\]
At coordinate \(n\), the ratio of the right side of
\eqref{eq:domination} to its left side is at most
\[
             Cq\exp(-tL_k+L_{k-1})\longrightarrow0.
\]
This contradicts \eqref{eq:domination}.
\end{proof}

\subsection{Idempotence, properness, and generator number}

Define
\[
 \mathcal G=
 \{D_{\xi_A^{\,2^{-r}}}:A\in\mathscr A,\ r\in\mathbb N_0\},
 \qquad
 \Istar=(\mathcal G).
\]
Every generator is compact, so \(\Istar\subset\KH\).

\begin{lemma}\label{lem:properties}
The ideal \(\Istar\) satisfies:
\begin{enumerate}[label=\textup{(\roman*)}]
 \item \(\Istar\subsetneq\KH\);
 \item \(\Istar^2=\Istar\);
 \item every \(\BH\)-generating set for \(\Istar\) has cardinality
       at least \(\cc\).
\end{enumerate}
\end{lemma}

\begin{proof}
For (i), let
\[
                    \zeta_n=\frac1{\log(n+1)}.
\]
This belongs to \(c_0^*\).  If \(D_\zeta\in\Istar\), finite ideal
membership and Lemma \ref{lem:finite-domination} would bound \(\zeta\) by
finitely many ampliations of \(\xi_{B_i}^{t_i}\).  Choose infinitely many
\(k\in\bigcap_iB_i\), which Lemma \ref{lem:independent} supplies, and again
evaluate at \(n=MN_k\).  With \(t=\min_i t_i\), the ratio of the proposed
right side to \(\zeta_n\) is at most
\[
          Cq\log(MN_k+1)e^{-tL_k}\longrightarrow0.
\]
Thus \(D_\zeta\notin\Istar\), proving strict containment.

For (ii), every generator has the exact factorization
\[
 D_{\xi_A^{\,2^{-r}}}
   =\left(D_{\xi_A^{\,2^{-(r+1)}}}\right)^2.
\]
Therefore \(\mathcal G\subset\Istar^2\).  Since \(\Istar^2\) is a
\(\BH\)-ideal containing \(\mathcal G\), minimality gives
\(\Istar\subset\Istar^2\).  The reverse inclusion is automatic for every
ideal.

For (iii), suppose that \(S\subset\Istar\), \(|S|<\cc\), and
\((S)=\Istar\).  Every element \(T\in S\) belongs to the ideal generated by
a finite subset of \(\mathcal G\).  Let \(U\subset\mathscr A\) be the union
of the corresponding finite sets of \(A\)-indices.  Then
\(|U|<\cc\); no regularity assumption on \(\cc\) is involved, since a union
of \(\kappa\) finite sets has cardinality at most \(\kappa\) for every
infinite cardinal \(\kappa\).

Choose \(A\in\mathscr A\setminus U\).  The ideal \((S)\) is contained in
the ideal generated by
\[
       \{D_{\xi_B^{2^{-r}}}:B\in U,\ r\in\mathbb N_0\}.
\]
Lemma \ref{lem:separation} says that this ideal does not contain
\(D_{\xi_A}\), contradicting \((S)=\Istar\).
\end{proof}

\begin{theorem}[Negative answer to Question 5]\label{thm:main}
There is a proper \(K(H)\)-ideal \(\mathcal I\subsetneq K(H)\) such that
\[
 \dim_{\mathbb C}(\mathcal I/\mathcal I^0)=0<\cc,
 \qquad
 \min\{|S|:(S)_{K(H)}=\mathcal I\}=\cc.
\]
Hence both assertions asked in Section 6, Question 5 of
arXiv:1209.6323 have negative answers in ZFC.
\end{theorem}

\begin{proof}
Take \(\mathcal I=\Istar\) and the ambient ideal \(J=\KH\).  If
\((S)_{\KH}=\Istar\), taking the \(\BH\)-ideal hull of both sides gives
\[
                         (S)=\Istar.
\]
Lemma \ref{lem:properties}(iii) therefore yields \(|S|\geq\cc\).

For the matching upper bound, idempotence gives
\(\Istar^3=\Istar\).  Since \(\Istar\subset\KH\),
\[
 \Istar=\Istar^3
 \subset\KH\,\Istar\,\KH
 =\KH(\mathcal G)\KH
 \subset(\mathcal G)_{\KH}
 \subset(\mathcal G)=\Istar.
\]
Thus \((\mathcal G)_{\KH}=\Istar\).  The family \(\mathcal G\) has
cardinality \(\cc\), so the minimum is exactly \(\cc\).

By Lemma \ref{lem:properties}(ii) and \(\Istar\subset\KH\),
\[
       \Istar=\Istar^2\subset\Istar\KH\subset\Istar.
\]
Thus \(\Istar\KH=\Istar\).  The defining expression for
\(\mathcal I^0\) contains \(\mathcal I J=\mathcal I\), while it is always
contained in \(\mathcal I\).  Hence
\(\mathcal I^0=\mathcal I\), and the quotient has dimension zero.
\end{proof}

\section{Verification notes}

The proof was subjected to the following checks.

\begin{enumerate}
 \item The independent family is constructed explicitly on a countable set;
       no additional set-theoretic axiom is used.
 \item Finite ideal membership is converted to ampliation domination using
       only a finite operator representation and the Ky Fan inequality.
 \item The choice \(n=MN_k\) handles all finitely many ampliations
       simultaneously and remains in the same block for sufficiently large
       \(k\).
 \item The separation lemma allows arbitrary positive powers, so adjoining
       all dyadic roots does not weaken the generator obstruction.
 \item The slowly decreasing \(\zeta\) proves
       \(\Istar\subsetneq\KH\); the counterexample is a proper subideal of
       its ambient ideal.
 \item The equality \(\Istar^2=\Istar\) is algebraic: each generator is an
       exact product.  No closure operation occurs.
 \item A set of fewer than \(\cc\) ideal elements involves fewer than
       \(\cc\) original indices because each element uses only finitely many
       algebraic generators.
 \item Taking \(\BH\)-ideal hulls of
       \((S)_{\KH}=\Istar\) gives \((S)=\Istar\), so the ordinary generator
       lower bound transfers to subideal generation.  Conversely,
       \(\Istar^3=\Istar\subset\KH\Istar\KH\) proves that the displayed
       continuum family really generates \(\Istar\) as a \(K(H)\)-ideal.
\end{enumerate}

The reusable script \texttt{code/verify\_block\_estimates.py} evaluates the
two logarithmic upper bounds for sample parameters.  By block \(k=4\) they
are already approximately \(-3.8\cdot10^3\) and
\(-4.1\cdot10^3\).  This is a scale-arithmetic consistency check only; it is
not part of the proof.

\section{Scope, limitations, and novelty check}

Theorem \ref{thm:main} answers Question 5 only.  It does not claim to settle
Questions 1--4.

The example has zero quotient dimension.  This is included in the literal
wording of Question 5, and \(\mathcal I\) is a proper subideal of the ambient
\(K(H)\).  If the intended question had the additional unstated hypothesis
\(\dim(\mathcal I/\mathcal I^0)>0\), that variant would remain open.

The bounded novelty search on 23 July 2026 covered:
\begin{itemize}
 \item the run's registry, solution, attempt, and proof-gap indexes;
 \item arXiv:1209.6323, its exact title, and exact phrases from Question 5;
 \item searches combining subideals, quotient dimension, generator
       cardinality, idempotent \(\BH\)-ideals, and continuum generators;
 \item the later Patnaik--Weiss survey arXiv:1303.5697 \cite{PWsurvey}.
\end{itemize}
The survey develops subideal traces and commutator spaces, but no later source
explicitly answering the generator-dimension Question 5 was located.  This
supports moderate novelty confidence, not an exhaustive claim of priority.

\textbf{Human review recommendation.}
Send the packet to an operator-ideal expert.  The decisive checks are Lemma
\ref{lem:finite-domination}, the finite-support argument in Lemma
\ref{lem:properties}(iii), and the source convention for \(J\)-generation.

\begin{thebibliography}{9}

\bibitem{PW12}
S. Patnaik and G. Weiss,
\emph{Subideals of operators II},
arXiv:1209.6323 (2012),
\url{https://arxiv.org/abs/1209.6323}.

\bibitem{Simon}
B. Simon,
\emph{Trace Ideals and Their Applications},
second edition, Mathematical Surveys and Monographs 120,
American Mathematical Society, 2005.

\bibitem{PWsurvey}
S. Patnaik and G. Weiss,
\emph{A survey on subideals of operators and an introduction to
subideal-traces},
arXiv:1303.5697 (2013),
\url{https://arxiv.org/abs/1303.5697}.

\end{thebibliography}

\end{document}
